\subsection{测度的窄收敛与函数的紧收敛}

命题 12. --- 设 T 是完全正则空间，B 是赋范空间 \(\mathcal{C}^{b}(\mathrm{T};\mathbf{C})\) 的单位球。设 I 是 \(\mathcal{C}^{b}(\mathrm{T};\mathbf{C})\) 上的线性形式。要使 T 上存在有界复测度 \(\theta\)，满足 \(\mathrm{I}(f)=\theta(f)\) 对所有 \(f\in\mathcal{C}^{b}(\mathrm{T};\mathbf{C})\) 成立，必要且充分的条件是 I 在 B 上关于紧收敛拓扑连续。测度 \(\theta\) 于是唯一。

证明陈述中的条件是必要的。设 \(\theta\) 是 \(\mathbf{T}\) 上的有界复测度，\(\varepsilon\) 是数 \(>0\)，\(\mathbf{K}\) 是 \(\mathbf{T}\) 的紧子集，满足 \(|\theta|^{\bullet}(\mathbf{T} - \mathbf{K}) < \varepsilon\)。取 \(f \in \mathbf{B}\)；记 \(\mathbf{U}\) 为 \(f\) 中 \(\mathbf{B}\) 关于紧收敛拓扑的邻域，它由满足 \(g \in \mathbf{B}\) 的函数 \(\sup_{x \in \mathbf{K}} |g(x) - f(x)| \leqslant \varepsilon\) 构成。于是，对每个 \(g \in \mathbf{U}\)，

\[
| \theta (g) - \theta (f) | \leqslant \int_ {\mathrm{T}} | g - f | d | \theta | \leqslant \varepsilon | \theta | ^ {\bullet} (\mathrm{K}) + 2 | \theta | ^ {\bullet} (\mathrm{T} - \mathrm{K}) \leqslant (\| \theta \| + 2) \varepsilon ,
\]

因为 \(|g - f|\) 的上界为 \(\varepsilon\)（在 \(K\) 上），而在 \(T - K\) 上的上界为 2。

反之，设线性形式 I 定义在 \(\mathcal{C}^{b}(\mathrm{T};\mathbf{C})\) 上，且其在 B 上关于紧收敛拓扑连续。于是，对每个数 \(\varepsilon > 0\)，存在数 a \textgreater{} 0 和 T 的紧子集 K，使得关系 \(f \in B\)、\(\sup_{x \in K} |f(x)| \leqslant a\) 蕴含 \(|\mathrm{I}(f)| \leqslant \varepsilon\)。第 2 号命题 5 遂蕴含存在唯一的有界复测度 \(\theta\)，使得 \(\mathrm{I}(f) = \theta(f)\) 对所有 \(f \in \mathcal{C}^{b}(\mathrm{T};\mathbf{C})\) 成立。

命题 13. --- 设 T 是局部紧空间，H 是赋范空间 \(\mathcal{C}^{b}(\mathrm{T};\mathbf{C})\) 的有界子集。映射 \((\mu,f)\mapsto\mu(f)\) 从 \(\mathcal{M}_{+}^{b}(\mathrm{T})\times\mathrm{H}\) 到 C 连续，只要 \(\mathcal{M}_{+}^{b}(\mathrm{T})\) 赋予窄拓扑、H 赋予紧收敛拓扑。

取 \(\mu \in \mathcal{M}_{+}^{b}(\mathrm{T})\)、\(f \in \mathrm{H}\)，取实数 \(\mathbf{M}\)，使得 \(\| \mu \| < \mathbf{M}\)，且对所有 \(|g| \leqslant \mathbf{M}\) 有 \(g \in \mathbf{H}\)。令 \(\varepsilon\) 为满足 \(>0\) 的数，取 \(\mathbf{K}\) 为 \(\mathbf{T}\) 的紧子集，使得 \(\mu^{\bullet}(\mathbf{T} - \mathbf{K}) < \varepsilon\)，再取 \(\mathbf{S}\) 为 \(\mathbf{K}\) 的相对紧开邻域。记 \(\mathbf{U}\) 为满足以下不等式的测度 \(\lambda \in \mathcal{M}_{+}^{b}(\mathbf{T})\) 所成的集合

\[
\lambda^ {\bullet} (\mathrm{T}) <   \mathrm{M}, \quad \lambda^ {\bullet} (\mathrm{T} - \mathrm{S}) <   \varepsilon , \quad | \lambda (f) - \mu (f) | <   \varepsilon
\]

于是，集合 U 是 \(\mu\) 在 \(\mathcal{M}_+^b (\mathrm{T})\) 中的邻域（第 3 号命题 6）。此外，令 V 为 H 中 \(f\) 的邻域，由函数 \(g\in \mathbf{H}\) 所成，且满足

\[
\sup _ {x \in S} | g (x) - f (x) | <   \varepsilon .
\]

取 \(\lambda \in \mathbf{U}\) 和 \(g\in \mathbf{V}\)；由于函数 \(|g - f|\) 在 \(\varepsilon\) 上被 \(\mathbf{S}\) 控制，在 \(\mathbf{T} - \mathbf{S}\) 上被 2M 控制，有

\[
\left| \lambda (g) - \lambda (f) \right| \leqslant \int_ {\mathrm{T}} | g - f | d \lambda \leqslant \varepsilon \lambda^ {\bullet} (\mathrm{S}) + 2 \mathrm{M} \lambda^ {\bullet} (\mathrm{T} - \mathrm{S}) \leqslant 3 \mathrm{M} \varepsilon ,
\]

由此得到

\[
\left| \lambda (g) - \mu (f) \right| \leqslant \left| \lambda (g) - \lambda (f) \right| + \left| \lambda (f) - \mu (f) \right| \leqslant (3 M + 1) \varepsilon .
\]

这就证明了映射 \((\lambda, g) \mapsto \lambda(g)\) 在 \((\mu, f)\) 这个 \(\mathcal{M}_+^b(\mathrm{T}) \times \mathrm{H}\) 中的点处连续。

备注。 --- 设 T 是完全正则空间，M 是满足 Prokhorov 条件的 \(\mathcal{M}^{b}(\mathrm{T})\) 的子集，H 是 \(\mathcal{C}^{b}(\mathrm{T})\) 的有界子集。可以用与刚才非常接近的论证证明映射 \((\lambda,g)\mapsto\lambda(g)\) 从 \(M\times H\) 到 C 连续，只要 M 赋予窄拓扑、H 赋予紧收敛拓扑。

推论。 --- 设 T 是完全正则空间，X 是拓扑空间，f 是定义在 \(T \times X\) 上的有界连续复值函数。对 T 上的每个有界测度 \(\mu\)，令 \(F_{\mu}\) 为 X 上由 \(\mathrm{F}_{\mu}(x) = \int_{\mathrm{T}} f(t, x) d\mu(t)\) 定义的函数，对所有 \(x \in X\) 成立。

\begin{enumerate}
\def\labelenumi{\alph{enumi})}
\item
  函数 \(\mathrm{F}_{\mu}\) 对每个有界测度 \(\mu\) 都连续且有界。
\item
  假设 T 局部紧。映射 \(\mu\mapsto F_{\mu}\) 从 \(\mathcal{M}_{+}^{b}(T)\) 到 \(\mathcal{C}^{b}(X;C)\) 连续，只要 \(\mathcal{M}_{+}^{b}(T)\) 赋予窄拓扑、\(\mathcal{C}^{b}(X;C)\) 赋予紧收敛拓扑。
\end{enumerate}

对每个 \(x \in X\)，记 \(f_{x}\) 为 T 上的连续有界函数 \(t \mapsto f(t, x)\)；映射 \(x \mapsto f_{x}\) 从 X 到 \(\mathcal{C}^{b}(\mathrm{T}; \mathbf{C})\) 的像有界，并且当 \(\mathcal{C}^{b}(\mathrm{T}; \mathbf{C})\) 赋予紧收敛拓扑时它连续（GT, X, §3, No.~4, Th. 3）。由于 \(\mathrm{F}_{\mu}(x) = \mu(f_{x})\)，由命题 12，函数 \(F_{\mu}\) 连续。假设 T 局部紧；命题 13 表明，映射 \((\mu, x) \mapsto \mathrm{F}_{\mu}(x)\) 从 \(\mathcal{M}_{+}^{b}(\mathrm{T}) \times \mathrm{X}\) 到 C 连续；断言 b) 由此得到（同上）。

